\documentclass{beamer}

\usetheme{CambridgeUS}
\usecolortheme{beaver}

\usepackage[utf8]{inputenc}
\usepackage[brazil]{babel}
\usepackage{caption}
% \usepackage{hyperref}
\hypersetup{
    final=true,
    % bookmarks=true,
    unicode=false,          % non-Latin characters in Acrobat’s bookmarks
    pdftoolbar=true,        % show Acrobat’s toolbar?
    pdfmenubar=true,        % show Acrobat’s menu?
    pdffitwindow=false,     % window fit to page when opened
    pdfstartview={FitH},    % fits the width of the page to the window
    pdftitle={Introdução ao \LaTeX com o Visual Studio Code no Ubuntu},    % title
    pdfauthor={Nightwind},     % author
    pdfsubject={Introdução ao Visual Studio Code no Ubuntu},   % subject of the document
    % pdfcreator={Creator},   % creator of the document
    % pdfproducer={Producer}, % producer of the document
    pdfkeywords={latex, dicas}, % list of keywords
    pdfnewwindow=true,      % links in new PDF window
    colorlinks=false,       % false: boxed links; true: colored links
    linkcolor=red,          % color of internal links (change box color with linkbordercolor)
    citecolor=green,        % color of links to bibliography
    filecolor=cyan,         % color of file links
    urlcolor=magenta        % color of external links
}
\usepackage[style=abnt]{biblatex}
\addbibresource{references.bib}
\usepackage{csquotes}

\usepackage{xcolor}
\usepackage{listings}
\lstset
{breakatwhitespace=true,breaklines=true,language=bash,basicstyle=\small\ttfamily,flexiblecolumns,backgroundcolor=\color{white}}
\usepackage{menukeys}

\title{Utilizando o Visual Studio Code no Ubuntu}
\author{Nightwind}
\institute[CTISM]{Colégio Técnico Industrial de Santa Maria}
\logo{\includegraphics[width=1cm]{../images/photo1.jpg}}
\date{\today}

\begin{document}

\frame{\titlepage}

\begin{frame}
    \frametitle{Sumário}
    \tableofcontents
\end{frame}

\section{Instalação do Visual Studio Code}

\begin{frame}[fragile]
    \frametitle{Instalação do Visual Studio Code}
    \begin{itemize}
        \item Abrir o terminal do Ubuntu: \keys{\ctrlwin + \Altwin + T }
    \end{itemize}
    \begin{lstlisting}
        $ sudo snap install --classic code # or code-insiders
    \end{lstlisting}
\end{frame}


\section{Instalação dos Requisitos}
\begin{frame}[fragile]
    \frametitle{Instalação dos Requisitos}
    \begin{lstlisting}
        $ sudo apt install texlive-full
    \end{lstlisting}
\end{frame}

\section{Referências}

\begin{frame}[allowframebreaks]
    \frametitle{Referências}
    \nocite{TeXLiveUbuntu,VisualStudioCodeUbuntu}
    \printbibliography[]
\end{frame}

\end{document}
